\documentclass[runningheads,oneside]{llncs}

% =====================
% Preamble (ONE place only)
% =====================
% =========================
% CORE FORMATTING
% =========================
\usepackage[utf8]{inputenc}
\usepackage[T1]{fontenc}
\usepackage{lmodern}
\usepackage{microtype}
\usepackage[margin=1in]{geometry}
\usepackage{parskip}
\usepackage{titlesec}

% =========================
% MATH
% =========================
\usepackage{amsmath, amssymb, amsthm}
\usepackage{mathtools}
\usepackage{bm} % bold math symbols

% =========================
% FIGURES & TABLES
% =========================
\usepackage{graphicx}
\usepackage{subcaption}
\usepackage{booktabs}
\usepackage{multirow}
\usepackage{array}
\usepackage{float} % better control of figure placement

% =========================
% COLOURS & CODE
% =========================
\usepackage{xcolor}
\usepackage{listings}

% Optional: nicer code formatting (safe defaults)
\lstset{
    basicstyle=\ttfamily\small,
    breaklines=true,
    frame=single,
    columns=fullflexible
}

% =========================
% LINKS
% =========================
\usepackage{hyperref}

\hypersetup{
    colorlinks = true,
    linkcolor  = black,
    citecolor  = black,
    urlcolor   = blue,
}

% =========================
% CITATIONS
% =========================
\usepackage[numbers]{natbib}

% =========================
% UTILITIES
% =========================
\hypersetup{
    colorlinks = true,
    linkcolor  = black,
    citecolor  = black,
    urlcolor   = blue,
}

\newcommand{\todo}[1]{\textcolor{red}{\textbf{[TODO: #1]}}}
\renewcommand{\thesubsection}{\thesection.\arabic{subsection}}
% \usepackage{showframe}


\title{Applied Deep Learning Report COMP0197 - Hierarchical Sales Forecasting Using Deep Learning}
\author{
Jordan Sydney-Darling \and
Imasha [surname] \and
Luke [surname]\and
Mia [surname]\and
Selin [surname]\and
Sherwin [surname]\and
Zhen-Yen [surname]\and
Ellen [surname]\and
Troy [surname]
}

\institute{
University College London, UK
}
\authorrunning{Sydney-Darling et al.}
\date{March 2026}

% =====================
\begin{document}
% =====================

\maketitle

\begin{abstract}
\todo{To be completed after results. Summarise objectives, models, and key findings.}
\end{abstract}

% =====================
% SECTIONS (modular)
% =====================

\section{Introduction}

\todo{Motivate probabilistic forecasting on M5 (uncertainty + hierarchy) and briefly state models compared (GRU/LSTM/LGBM/TFT).}

Our implementation, including all models and experiments, is publicly available\footnote{Code: \url{https://github.com/...}}.
\section{Data and Preprocessing}

\subsection{Dataset}
\todo{Describe M5 dataset (hierarchical structure, 28-day horizon, sparsity/intermittency).}

\subsection{Feature Engineering}
\todo{Summarise key features (lags, rolling stats, price, calendar/events) and why they matter for demand.}

\subsection{Data Pipeline}
\todo{Outline data flow (split, normalisation, batching) shared across all models.}

\section{Methods}
\subsection{Baselines}

\subsubsection{GRU Baseline Ablation Study}
This section presents a controlled ablation study on GRU-based models, including both methodological setup and corresponding results, to isolate the effect of predictive objective under fixed architectural conditions.\\\\
We implement five GRU-based baseline variants to isolate the effect of predictive objective on forecasting performance while controlling for model architecture. All models share a common backbone consisting of a single GRU layer with the final hidden state passed to a linear output head, ensuring that differences in performance arise from the loss formulation rather than architectural complexity. The variants include: a deterministic baseline trained with MSE, a Gaussian probabilistic model outputting $(\mu, \sigma)$ trained via negative log-likelihood, a Negative Binomial model parameterised by $(\mu, \alpha)$ to account for overdispersed count data, a quantile regression model trained using pinball loss across seven quantiles $[0.025, 0.05, 0.25, 0.5, 0.75, 0.95, 0.975]$, and a revenue-weighted quantile model where per-series losses are scaled by volume $\times$ average sell price. 

To further disentangle modelling choices, we conducted controlled experiments to evaluate the impact of forecasting strategy and feature representation. Each model was trained under identical settings while varying only the decoding regime (autoregressive versus direct multi-step/seq2seq) and the treatment of hierarchical information (flat encoded features versus learned hierarchical representations). Hyperparameters were selected using successive halving on a stratified subset of 3,000 series, with final models trained on the full dataset of 30,490 series. This setup allows us to isolate the contribution of loss function, forecasting strategy, and hierarchy modelling independently of architectural changes.


\paragraph{Results and Discussion}

\todo{2--3 sentences. State the deterministic baseline RMSE as the lower
bound. State which probabilistic variant performed best on point accuracy and
whether it matched or beat the baseline. One sentence on best calibration
result (coverage closest to 0.95).}

\begin{table}[h]
\centering
\caption{\textbf{GRU baseline comparison across forecasting regimes and hierarchy representations on the M5 test set (d1914--d1941, 30,490 series).} Best per column in \textbf{bold}. AR = autoregressive; S2S = direct multi-step.}
\label{tab:gru_results}

\renewcommand{\arraystretch}{1.5} 

\resizebox{\textwidth}{!}{ 
\begin{tabular}{lllrrrrrrrr}
\toprule
\textbf{Model} & \textbf{Regime} & \textbf{Hierarchy}
    & \textbf{RMSE} & \textbf{W-RMSE}
    & \textbf{MAE}  & \textbf{W-MAE}
    & \textbf{$R^2$} & \textbf{Cov@95\%} & \textbf{Width} \\
\midrule

Deterministic (MSE)  & AR  & Flat 
& 2.74 & 0.01781 
& 1.14 & 0.000045 
& 0.434 & -- & -- \\

Gaussian (NLL)       & AR  & Flat 
& 2.78 & 0.01777 
& 1.20 & 0.000046 
& 0.419 & 0.824 & 2.40 \\

Quantile (Pinball)   & AR  & Flat 
& \textbf{2.29} & \textbf{0.01518} 
& \textbf{0.97} & \textbf{0.000039} 
& \textbf{0.605} & \textbf{0.892} & 3.36 \\
\midrule

Deterministic (MSE) & S2S & Flat & XX.XX & XX.XX & XX.XX & XX.XX & X.XXX & XX.X & XX.X \\

Best Parametric & S2S & Flat & XX.XX & XX.XX & XX.XX & XX.XX & X.XXX & XX.X & XX.X \\

Best Non-Parametric & S2S & Flat & XX.XX & XX.XX & XX.XX & XX.XX & X.XXX & XX.X & XX.X \\
\midrule

Hierarchical Deterministic & S2S & Hierarchical & XX.XX & XX.XX & XX.XX & XX.XX & X.XXX & XX.X & XX.X \\

Best Hierarchical Parametric & S2S & Hierarchical & XX.XX & XX.XX & XX.XX & XX.XX & X.XXX & XX.X & XX.X \\

Best Hierarchical Non-Parametric & S2S & Hierarchical & XX.XX & XX.XX & XX.XX & XX.XX & X.XXX & XX.X & XX.X \\
\midrule

Best Hierarchical & AR  & Hierarchical & XX.XX & XX.XX & XX.XX & XX.XX & X.XXX & XX.X & XX.X \\

\bottomrule
\end{tabular}
}

\end{table}

Direct multi-step (seq2seq) forecasting consistently outperformed autoregressive training across flat baselines, indicating that predicting the full horizon jointly is more effective for this task. Incorporating hierarchical embeddings further improved performance, suggesting that explicit modelling of product hierarchy is beneficial beyond using raw categorical features.

Based on these findings, we recommend using direct multi-step forecasting with explicit hierarchy modelling for downstream models such as TFT.

However, while GRU-based models capture local temporal dynamics effectively, they lack the ability to model long-range dependencies and heterogeneous feature interactions, motivating the use of more expressive architectures such as TFT.          
\subsubsection{LSTM Baseline Study}

\todo{Mirror GRU setup (same inputs + losses) and briefly contrast behaviour vs GRU.}

\subsection{LightGBM Quantile Regression}

\todo{Explain tabular approach (lags/features → trees) and quantile prediction via pinball loss.}
\subsection{Hierarchical LSTM}

\todo{Describe use of hierarchy (state/store/category) to improve sparse series via shared patterns.}
\subsection{Temporal Fusion Transformer}

\todo{Explain multi-horizon attention model using static + temporal features with direct quantile outputs.}

\section{Experiments}

\subsection{Training Procedure}
\todo{Summarise shared training setup (optimiser, epochs, early stopping, same conditions across models).}

\subsubsection*{Pinball vs Revenue Weighted Pinball Loss}

Quantile regression is trained using the pinball loss, which estimates conditional quantiles \( \tau \in (0,1) \):

\begin{equation}
\mathcal{L}_{\tau}(\hat{y}, y) =
\begin{cases}
\tau (y - \hat{y}) & \text{if } y \geq \hat{y} \\
(1 - \tau)(\hat{y} - y) & \text{otherwise}
\end{cases}
\end{equation}

This asymmetric loss enables direct learning of prediction intervals without assuming a parametric distribution.

However, standard pinball loss treats all items equally, which is unsuitable for retail forecasting. Forecast errors have unequal business costs: underestimating high-revenue products leads to unmet demand and lost sales, while overestimation leads to excess inventory. These costs scale with price and volume.

We therefore use a revenue-weighted variant:

\begin{equation}
\mathcal{L}_{\text{weighted}} =
\sum_{i=1}^{N} w_i \sum_{\tau \in \mathcal{T}} \mathcal{L}_{\tau}(\hat{y}_{i,\tau}, y_i),
\quad \sum_i w_i = 1
\end{equation}

where \( w_i \propto \) (sales volume \( \times \) average price).

\paragraph{} Equal weighting optimises average error; revenue weighting optimises financial impact.

\subsection{Evaluation Metrics}
\todo{Define metrics (RMSE, MAE, weighted variants, pinball, coverage, interval width).}

\section{Results}

\subsection{Baseline Comparison}
\todo{Compare GRU/LSTM/LGBM on point accuracy and identify strongest baseline.}

\subsection{Advanced Model Performance}
\todo{Compare hierarchical LSTM and TFT vs baselines (do they improve accuracy?).}

\subsection{Uncertainty Evaluation}
\todo{Evaluate calibration (coverage vs width) across probabilistic models.}
\section{Discussion}

\todo{Interpret key findings (does probabilistic + hierarchy + TFT help?) and note limitations.}
\section{Conclusion}

\todo{Summarise main takeaway (best model + why) and overall insight from comparison.}

\bibliographystyle{splncs04}
\bibliography{references}


% =====================
\newpage
\appendix
\section{Notes on Writing and Compilation}

The report has already been split into individual \texttt{.tex} files, one per section. The full \texttt{Report} folder (including all section files and a \texttt{main.tex}) has been uploaded to Git and to a shared Overleaf project.

The recommended approach is to work directly in Overleaf. Each person should open the project, locate their assigned \texttt{.tex} file, and write directly within it. This allows everyone to see updates in real time and ensures the full report always compiles in one place.

Each file already contains the appropriate section heading. There is no need to include a preamble or document setup — only the content for that section should be written.

For example, a file might look like:

\begin{verbatim}
\subsection{LSTM Baseline Study}

Write your content here...
\end{verbatim}

Nothing else is required in the file (no \texttt{\textbackslash begin\{document\}}, \texttt{\textbackslash end\{document\}}, or additional setup).

The \texttt{main.tex} file compiles all sections into a single document using \texttt{\textbackslash input\{\}}.

\paragraph{}

A shared \texttt{preamble.tex} file contains all packages and settings for the report and is already imported into \texttt{main.tex}. There is no need to include packages in individual section files. If anything additional is needed, it can be added to the end of \texttt{preamble.tex}.

A shared \texttt{references.bib} file is included for citations. If you add a reference, include it there and cite it using \texttt{\textbackslash cite\{\}} in your section.

An example reference \cite{makridakis2022m5}

\paragraph{}

If preferred, sections can also be edited locally and then uploaded, but working within Overleaf is recommended to keep everything synchronised.


\begin{center}
\fbox{
\parbox{0.9\textwidth}{
\centering
\textbf{\Large Overleaf Project Link}\\[6pt]
\texttt{https://www.overleaf.com/6558942236hxyjdxffjntn#b4b1cf}
}
}
\end{center}

\newpage

\section{USEFUL STUFF}
USEFUL EQUATIONS WE MIGHT DECIDE TO PUT IN THE REPORT
\subsubsection*{Probabilistic Gaussian}
\begin{equation}
    \mathcal{L}_{\mathrm{Gauss}}
    \;=\;
    \frac{1}{N} \sum_{i=1}^{N}
    \left[
        \log \sigma_i
        + \frac{(y_i - \mu_i)^2}{2\,\sigma_i^2}
    \right]
    + \lambda\,\bar{\sigma}
    \label{eq:gauss_nll}
\end{equation}


\subsubsection*{Probabilistic Negative Binomial}

The Negative Binomial log-likelihood is:
%
\begin{equation}
    \mathcal{L}_{\mathrm{NB}}
    \;=\;
    -\frac{1}{N} \sum_{i=1}^{N}
    \Bigl[
        \log\Gamma(y_i + r_i)
        - \log\Gamma(r_i)
        - \log\Gamma(y_i + 1)
        + r_i \log\!\frac{r_i}{r_i + \mu_i}
        + y_i \log\!\frac{\mu_i}{r_i + \mu_i}
    \Bigr]
    \label{eq:nb_nll}
\end{equation}
%
where $r_i = 1/\alpha_i$ is the dispersion parameter.

\subsubsection*{Quantile Regression}

The pinball loss for quantile $\tau$ is:
%
\begin{equation}
    \mathcal{L}_{\tau}(\hat{y},\, y)
    \;=\;
    \begin{cases}
        \tau\,(y - \hat{y})       & \text{if } y \geq \hat{y} \\[4pt]
        (1 - \tau)\,(\hat{y} - y) & \text{otherwise}
    \end{cases}
    \label{eq:pinball}
\end{equation}



\subsubsection*{Forecast Visualisations}
COPY TEMPLATE IF ANYONE WANTS TO PUT VISUALSE IN

\begin{figure}[h]
    \centering
    \begin{subfigure}[b]{0.32\textwidth}
        \centering
        \includegraphics[width=\textwidth]{forecast_1.png}
        \caption{Forecast series 1}
    \end{subfigure}
    \hfill
    \begin{subfigure}[b]{0.32\textwidth}
        \centering
        \includegraphics[width=\textwidth]{forecast_2.png}
        \caption{Forecast series 2}
    \end{subfigure}
    \hfill
    \begin{subfigure}[b]{0.32\textwidth}
        \centering
        \includegraphics[width=\textwidth]{forecast_3.png}
        \caption{Forecast series 3}
    \end{subfigure}
    \caption{ADD A CAPTION}
    \label{fig:forecasts}
\end{figure}

\newpage


%=====================
\end{document}